\documentclass[aspectratio=169]{beamer}
\usetheme{Madrid}
\usecolortheme{seahorse}

\usepackage[italian]{babel}
\usepackage[utf8]{inputenc}
\usepackage[T1]{fontenc}
\usepackage{listings}
\usepackage{xcolor}
\usepackage{booktabs}
\usepackage{tikz}

% Colori personalizzati
\definecolor{sqlblue}{RGB}{0,82,155}
\definecolor{sqlgreen}{RGB}{0,128,0}
\definecolor{sqlred}{RGB}{180,0,0}
\definecolor{codefg}{RGB}{40,40,40}
\definecolor{codebg}{RGB}{248,248,248}
\definecolor{keyw}{RGB}{0,0,200}

% Stile listing per SQL
\lstdefinestyle{sqlstyle}{
  language=SQL,
  basicstyle=\ttfamily\footnotesize,
  keywordstyle=\color{keyw}\bfseries,
  commentstyle=\color{sqlgreen}\itshape,
  stringstyle=\color{sqlred},
  backgroundcolor=\color{codebg},
  frame=single,
  rulecolor=\color{gray!40},
  breaklines=true,
  showstringspaces=false,
  tabsize=2,
  numbers=none,
  xleftmargin=4pt,
  xrightmargin=4pt,
}

% Stile listing per output/risultati
\lstdefinestyle{outputstyle}{
  basicstyle=\ttfamily\tiny,
  backgroundcolor=\color{black!90},
  frame=single,
  rulecolor=\color{gray},
  breaklines=true,
  showstringspaces=false,
  xleftmargin=4pt,
  xrightmargin=4pt,
  basicstyle=\ttfamily\tiny\color{green!70!black},
}

\setbeamertemplate{navigation symbols}{}
\setbeamerfont{frametitle}{size=\large,series=\bfseries}

\title[Trigger SQL]{Trigger nelle Basi di Dati}
\subtitle{Paradigma E-C-A, Tipi, Sintassi MySQL ed Esempi}
\author{IIS Fermi Sacconi Ceci}
\institute{Prof. Fedeli Massimo -- Basi di Dati}
\date{\today}

\begin{document}

% =====================================================================
% SLIDE 1 - Titolo
% =====================================================================
\begin{frame}
  \titlepage
\end{frame}

% =====================================================================
% SLIDE 2 - Indice
% =====================================================================
\begin{frame}{Indice}
  \tableofcontents
\end{frame}

% =====================================================================
\section{Introduzione}
% =====================================================================

% SLIDE 3
\begin{frame}{Introduzione ai Trigger}
  \begin{block}{Definizione}
    L'introduzione di \textbf{trigger} all'interno di una Base di Dati permette
    la \textbf{gestione automatica} di particolari procedure in risposta a
    determinati eventi esterni.
  \end{block}

  \vspace{0.4cm}
  \begin{alertblock}{Basi di Dati Attive}
    Basi di Dati che utilizzano trigger vengono dette \textbf{Basi di Dati attive}:
    i trigger rendono \emph{reattivo} il comportamento del sistema alle
    sollecitazioni esterne.
  \end{alertblock}

  \vspace{0.3cm}
  \begin{tikzpicture}
    \node[draw,rounded corners,fill=sqlblue!15,minimum width=2.5cm,minimum height=0.7cm] (ev) at (0,0) {Evento};
    \node[draw,rounded corners,fill=orange!20,minimum width=2.5cm,minimum height=0.7cm] (tr) at (4,0) {Trigger};
    \node[draw,rounded corners,fill=sqlgreen!20,minimum width=2.5cm,minimum height=0.7cm] (az) at (8,0) {Azione};
    \draw[->,thick] (ev) -- (tr);
    \draw[->,thick] (tr) -- (az);
  \end{tikzpicture}
\end{frame}

% SLIDE 4
\begin{frame}{Il Livello di Astrazione dei Dati}
  \begin{columns}[T]
    \column{0.55\textwidth}
      L'utilizzo di \textbf{regole attive} in un DB aumenta il livello di
      astrazione dei dati, introducendo una nuova dimensione:
      \begin{itemize}
        \item \textbf{Indipendenza della conoscenza}: nasconde all'esterno
              l'esecuzione delle regole attive.
      \end{itemize}
      \vspace{0.3cm}
      Il vantaggio principale è lo \textbf{spostamento} di alcuni controlli e
      procedure automatiche \textbf{dal livello applicativo} a quello dello schema.
    \column{0.42\textwidth}
      \begin{tikzpicture}
        \node[draw,fill=sqlblue!20,minimum width=3cm,minimum height=0.6cm,rounded corners]
          at (0,2.4) {\footnotesize Livello Applicativo};
        \node[draw,fill=orange!20,minimum width=3cm,minimum height=0.6cm,rounded corners]
          at (0,1.6) {\footnotesize Regole Attive (Trigger)};
        \node[draw,fill=sqlgreen!20,minimum width=3cm,minimum height=0.6cm,rounded corners]
          at (0,0.8) {\footnotesize Schema DB};
        \node[draw,fill=gray!20,minimum width=3cm,minimum height=0.6cm,rounded corners]
          at (0,0) {\footnotesize Dati Fisici};
        \draw[<->,thick,sqlblue] (-1.8,0.8) -- node[left]{\tiny sposta} (-1.8,1.6);
      \end{tikzpicture}
  \end{columns}
\end{frame}

% =====================================================================
\section{Il Paradigma E-C-A}
% =====================================================================

% SLIDE 5
\begin{frame}{Il Paradigma E-C-A}
  \begin{block}{}
    Il funzionamento dei trigger segue il paradigma \textbf{Evento -- Condizione -- Azione}
  \end{block}
  \vspace{0.2cm}
  \begin{columns}[T]
    \column{0.33\textwidth}
      \begin{exampleblock}{\textcolor{sqlred}{Evento}}
        Qualsiasi modifica DML (\texttt{INSERT}, \texttt{UPDATE}, \texttt{DELETE})
        o l'attivazione di un altro trigger (\emph{trigger in cascata}) su una
        \emph{tabella di target}.
      \end{exampleblock}
    \column{0.33\textwidth}
      \begin{exampleblock}{\textcolor{sqlblue}{Condizione} \small[opz.]}
        Un predicato booleano espresso mediante sintassi SQL. Se non verificato,
        il trigger non esegue l'azione.
      \end{exampleblock}
    \column{0.33\textwidth}
      \begin{exampleblock}{\textcolor{sqlgreen}{Azione}}
        Sequenza di primitive SQL, eventualmente arricchite da un linguaggio di
        programmazione integrato nel DBMS (es.\ PL/SQL, T-SQL).
      \end{exampleblock}
  \end{columns}
\end{frame}

% SLIDE 6
\begin{frame}{Trigger in Cascata}
  \begin{alertblock}{Trigger in cascata}
    Quando l'evento che attiva un trigger è a sua volta l'azione di un altro
    trigger, si parla di \textbf{trigger in cascata}.
  \end{alertblock}
  \vspace{0.4cm}
  \begin{tikzpicture}
    \node[draw,fill=sqlblue!15,rounded corners,minimum width=2.2cm,minimum height=0.65cm]
      (op) at (0,0) {\footnotesize Op.\ DML};
    \node[draw,fill=orange!20,rounded corners,minimum width=2.2cm,minimum height=0.65cm]
      (t1) at (3.5,0) {\footnotesize Trigger 1};
    \node[draw,fill=orange!20,rounded corners,minimum width=2.2cm,minimum height=0.65cm]
      (t2) at (7,0) {\footnotesize Trigger 2};
    \node[draw,fill=sqlgreen!20,rounded corners,minimum width=2.2cm,minimum height=0.65cm]
      (az) at (10.5,0) {\footnotesize Azione finale};
    \draw[->,thick] (op) -- (t1);
    \draw[->,thick] (t1) -- (t2) node[midway,above]{\tiny evento};
    \draw[->,thick] (t2) -- (az);
  \end{tikzpicture}
  \vspace{0.3cm}
  \begin{block}{Attenzione}
    I trigger in cascata devono essere progettati con attenzione per evitare
    cicli infiniti di attivazione.
  \end{block}
\end{frame}

% =====================================================================
\section{Tipi di Trigger}
% =====================================================================

% SLIDE 7
\begin{frame}{Tipi di Trigger -- Panoramica}
  \begin{center}
    \begin{tabular}{lll}
      \toprule
      \textbf{Classificazione} & \textbf{Tipo} & \textbf{Descrizione breve} \\
      \midrule
      Granularità & Row-level & Eseguito per ogni riga coinvolta \\
                  & Statement-level & Eseguito una sola volta per istruzione \\
      \midrule
      Momento     & BEFORE & Prima dell'evento DML \\
                  & AFTER  & Dopo l'evento DML \\
                  & INSTEAD OF & Al posto dell'evento (Oracle) \\
      \midrule
      Livello     & Tabella & Risponde a eventi DML su tabelle \\
                  & Schema/DB & Risponde a eventi DDL \\
      \bottomrule
    \end{tabular}
  \end{center}
\end{frame}

% SLIDE 8
\begin{frame}{Trigger a Livello di Tupla (Row-Level)}
  \begin{block}{Definizione}
    I trigger \textbf{a livello di riga} (row-level) vengono eseguiti
    \textbf{una volta per ogni riga} su cui agisce un'istruzione DML.
  \end{block}
  \vspace{0.3cm}
  \begin{exampleblock}{Esempio}
    Un'istruzione \texttt{UPDATE} che modifica 500 righe attiva il trigger
    \textbf{500 volte} -- una per ciascuna riga modificata.
  \end{exampleblock}
  \vspace{0.3cm}
  \begin{columns}[T]
    \column{0.5\textwidth}
      \textbf{Uso tipico:}
      \begin{itemize}
        \item Validazione per singola riga
        \item Calcoli derivati riga per riga
        \item Sincronizzazione dati distribuiti
        \item Audit trail dettagliato
      \end{itemize}
    \column{0.5\textwidth}
      \textbf{Parola chiave MySQL:}
      \begin{center}
        \texttt{\textbf{FOR EACH ROW}}
      \end{center}
      Obbligatoria in MySQL -- solo row-level è supportato.
  \end{columns}
\end{frame}

% SLIDE 9
\begin{frame}{Trigger a Livello di Istruzione (Statement-Level)}
  \begin{block}{Definizione}
    I trigger \textbf{a livello di istruzione} vengono eseguiti
    \textbf{una sola volta} per ogni istruzione DML, indipendentemente
    dal numero di righe coinvolte.
  \end{block}
  \vspace{0.3cm}
  \begin{exampleblock}{Esempio}
    Se una singola istruzione \texttt{INSERT} inserisce 500 righe, un trigger
    a livello di istruzione viene eseguito \textbf{una sola volta}.
  \end{exampleblock}
  \vspace{0.3cm}
  \begin{alertblock}{Nota}
    I trigger a livello di istruzione \textbf{non vengono utilizzati di frequente}
    per attività correlate ai singoli dati. Sono più adatti per operazioni di
    logging a livello globale di transazione.
    \\ \textbf{MySQL non supporta} statement-level trigger (\texttt{FOR EACH ROW} è obbligatorio).
  \end{alertblock}
\end{frame}

% SLIDE 10
\begin{frame}{Trigger BEFORE e AFTER}
  \begin{columns}[T]
    \column{0.5\textwidth}
      \begin{block}{BEFORE Trigger}
        \begin{itemize}
          \item Eseguito \textbf{prima} dell'operazione DML
          \item Può \textbf{modificare} i valori \texttt{NEW} prima che vengano scritti
          \item Se fallisce, l'operazione DML \textbf{non viene eseguita}
          \item Utile per: validazione, normalizzazione dati
        \end{itemize}
      \end{block}
    \column{0.5\textwidth}
      \begin{alertblock}{AFTER Trigger}
        \begin{itemize}
          \item Eseguito \textbf{dopo} l'operazione DML
          \item I dati sono già stati modificati
          \item Viene eseguito solo se BEFORE e DML hanno avuto successo
          \item Utile per: audit, aggiornamenti correlati, notifiche
        \end{itemize}
      \end{alertblock}
  \end{columns}
  \vspace{0.5cm}
  \begin{center}
    \begin{tikzpicture}
      \node[draw,fill=yellow!20,rounded corners] (b) at (0,0) {BEFORE};
      \node[draw,fill=sqlblue!20,rounded corners] (d) at (3.5,0) {DML};
      \node[draw,fill=sqlgreen!20,rounded corners] (a) at (7,0) {AFTER};
      \draw[->,thick] (b) -- (d);
      \draw[->,thick] (d) -- (a);
      \node[above=0.15cm,font=\tiny] at (0,0) {validazione};
      \node[above=0.15cm,font=\tiny] at (7,0) {azioni collaterali};
    \end{tikzpicture}
  \end{center}
\end{frame}

% SLIDE 11
\begin{frame}{Trigger INSTEAD OF}
  \begin{block}{Definizione}
    In alcuni DBMS (ad es.\ \textbf{Oracle}) è possibile utilizzare trigger
    \texttt{INSTEAD OF} per indicare cosa fare \textbf{al posto} di eseguire
    le azioni che hanno invocato il trigger.
  \end{block}
  \vspace{0.3cm}
  \begin{exampleblock}{Caso d'uso principale: aggiornamento di viste}
    Con un trigger \texttt{INSTEAD OF} è possibile specificare al DBMS come
    eseguire \texttt{UPDATE}, \texttt{DELETE} e \texttt{INSERT} nelle
    \textbf{tabelle base della vista} quando un utente tenta di modificare
    i valori attraverso la vista stessa.
  \end{exampleblock}
  \vspace{0.2cm}
  \begin{alertblock}{Supporto in MySQL}
    MySQL \textbf{non supporta} i trigger \texttt{INSTEAD OF}.
    Sono disponibili in Oracle, SQL Server e PostgreSQL.
  \end{alertblock}
\end{frame}

% SLIDE 12
\begin{frame}{Trigger a Livello di Schema e Database}
  \begin{block}{Definizione}
    Trigger che si innescano su operazioni eseguite a \textbf{livello di schema}:
  \end{block}
  \vspace{0.2cm}
  \begin{columns}[T]
    \column{0.5\textwidth}
      \textbf{Esempi di eventi DDL:}
      \begin{itemize}
        \item \texttt{CREATE TABLE}
        \item \texttt{ALTER TABLE}
        \item \texttt{DROP TABLE}
        \item \texttt{RENAME}
        \item \texttt{TRUNCATE}
        \item \texttt{REVOKE} / \texttt{AUDIT}
      \end{itemize}
    \column{0.5\textwidth}
      \textbf{Capacità principali:}
      \begin{enumerate}
        \item \textbf{Impedire} l'esecuzione di operazioni DDL non autorizzate
        \item \textbf{Controllo della sicurezza} aggiuntivo quando si verificano operazioni DDL
        \item \textbf{Logging} delle modifiche strutturali
      \end{enumerate}
  \end{columns}
\end{frame}

% =====================================================================
\section{I Trigger SQL:2003}
% =====================================================================

% SLIDE 13
\begin{frame}{I Trigger nello Standard SQL:2003}
  \begin{block}{Definizione formale}
    Un trigger è una \textbf{specifica per una determinata azione} che deve
    essere eseguita ogniqualvolta avviene un'azione su un determinato oggetto.
  \end{block}
  \vspace{0.2cm}
  \textbf{Terminologia dello standard:}
  \begin{description}
    \item[\textit{triggered action}] L'azione d'innesco: una dichiarazione di
      procedura SQL o una lista di dichiarazioni.
    \item[\textit{target table}] La tabella della base di dati persistente su cui
      il trigger è definito (tabella obiettivo).
    \item[\textit{event trigger}] L'evento scatenante: inserimento, cancellazione
      o modifica di un insieme di righe.
  \end{description}
  \vspace{0.2cm}
  \begin{center}
    \begin{tabular}{ll}
      \toprule
      \textbf{Per momento} & AFTER TRIGGER, BEFORE TRIGGER \\
      \textbf{Per evento}  & DELETE, INSERT, UPDATE TRIGGER \\
      \bottomrule
    \end{tabular}
  \end{center}
\end{frame}

% SLIDE 14
\begin{frame}{Tabelle di Transizione}
  \begin{block}{Scopo}
    Per ogni trigger possono essere usate \textbf{tabelle di transizione}
    per memorizzare lo stato \emph{precedente} e \emph{successivo} del target.
  \end{block}
  \vspace{0.3cm}
  \begin{columns}[T]
    \column{0.33\textwidth}
      \begin{exampleblock}{DELETE Trigger}
        Solo \textbf{old transition table}: contiene le righe prima
        della cancellazione.
      \end{exampleblock}
    \column{0.33\textwidth}
      \begin{exampleblock}{INSERT Trigger}
        Solo \textbf{new transition table}: contiene le righe appena
        inserite.
      \end{exampleblock}
    \column{0.33\textwidth}
      \begin{exampleblock}{UPDATE Trigger}
        \textbf{Entrambe} le tabelle: old (prima) e new (dopo la modifica).
      \end{exampleblock}
  \end{columns}
  \vspace{0.3cm}
  \begin{center}
    \renewcommand{\arraystretch}{1.3}
    \begin{tabular}{lcc}
      \toprule
      \textbf{Evento} & \textbf{OLD} & \textbf{NEW} \\
      \midrule
      INSERT & \texttimes & \checkmark \\
      DELETE & \checkmark & \texttimes \\
      UPDATE & \checkmark & \checkmark \\
      \bottomrule
    \end{tabular}
  \end{center}
\end{frame}

% SLIDE 15
\begin{frame}{La Granularità dell'Azione}
  L'azione del trigger può essere di tipo:
  \vspace{0.3cm}
  \begin{columns}[T]
    \column{0.5\textwidth}
      \begin{block}{Statement-Level Trigger}
        Eseguito \textbf{solo una volta} dopo l'attivazione del trigger,
        indipendentemente dal numero di righe coinvolte.
        \vspace{0.2cm}
        \begin{itemize}
          \item La variabile di transizione ha il valore dell'intera tabella
        \end{itemize}
      \end{block}
    \column{0.5\textwidth}
      \begin{alertblock}{Row-Level Trigger}
        Eseguito \textbf{una volta per ogni riga} della tabella di transizione
        associata al trigger.
        \vspace{0.2cm}
        \begin{itemize}
          \item La variabile di transizione assume il valore della singola tupla
        \end{itemize}
      \end{alertblock}
  \end{columns}
\end{frame}

% SLIDE 16
\begin{frame}{Variabili dei Trigger}
  Vengono messe a disposizione dei trigger variabili speciali per accedere alle
  tabelle di transizione.
  \vspace{0.3cm}
  \begin{columns}[T]
    \column{0.5\textwidth}
      \begin{block}{Statement-Level Trigger}
        La variabile ha il valore dell'intera \textbf{tabella di transizione}.
      \end{block}
      \vspace{0.3cm}
      \begin{block}{Row-Level Trigger}
        La \textbf{variabile di transizione} associa ad ogni azione del trigger
        il valore della tupla nella tabella di transizione da modificare.
      \end{block}
    \column{0.5\textwidth}
      \begin{exampleblock}{Tipi di variabile (row-level)}
        \begin{itemize}
          \item \textbf{old transition variable}: valore della tupla \emph{prima}
            della modifica (DELETE, UPDATE)
          \item \textbf{new transition variable}: valore della tupla \emph{dopo}
            la modifica (INSERT, UPDATE)
        \end{itemize}
      \end{exampleblock}
  \end{columns}
\end{frame}

% SLIDE 17
\begin{frame}{Definizione Completa di un Trigger}
  Un trigger è definito dai seguenti componenti:
  \vspace{0.2cm}
  \begin{columns}[T]
    \column{0.5\textwidth}
      \begin{enumerate}
        \item \textbf{Nome} del trigger
        \item \textbf{Tabella di target}
        \item \textbf{Momento d'azione}: \texttt{BEFORE} / \texttt{AFTER}
        \item \textbf{Evento}: \texttt{INSERT} / \texttt{UPDATE} / \texttt{DELETE}
      \end{enumerate}
    \column{0.5\textwidth}
      \begin{enumerate}\setcounter{enumi}{4}
        \item \textbf{Granularità}: statement-level / row-level
        \item \textbf{Nomi} per le tabelle di transizione o variabili
        \item \textbf{L'azione} SQL da eseguire
        \item \textbf{Timestamp} della creazione
      \end{enumerate}
  \end{columns}
  \vspace{0.4cm}
  \begin{exampleblock}{Schema sintetico}
    \texttt{CREATE TRIGGER} \textit{nome}
    \texttt{BEFORE|AFTER INSERT|UPDATE|DELETE ON} \textit{tabella}
    \texttt{FOR EACH ROW} \textit{azione}
  \end{exampleblock}
\end{frame}

% =====================================================================
\section{Trigger in MySQL}
% =====================================================================

% SLIDE 18
\begin{frame}{Contesto di Esecuzione del Trigger (TEC)}
  \begin{block}{Esecuzione}
    Durante l'esecuzione di un'istruzione S possono esistere \textbf{zero o più}
    contesti d'esecuzione trigger (TEC), ma \textbf{al più uno} di loro può
    essere attivo alla volta.
  \end{block}
  \vspace{0.3cm}
  \begin{block}{Ciclo di vita del TEC}
    Il contesto di esecuzione di un trigger (TEC) viene:
    \begin{enumerate}
      \item \textbf{Creato} al momento dell'evento scatenante
      \item \textbf{Attivato} se la condizione è verificata
      \item \textbf{Eseguito}: viene eseguita l'azione definita
      \item \textbf{Distrutto} al termine dell'azione
    \end{enumerate}
  \end{block}
\end{frame}

% SLIDE 19
\begin{frame}[fragile]{Sintassi CREATE TRIGGER in MySQL}
  \begin{lstlisting}[style=sqlstyle]
CREATE
  [DEFINER = { user | CURRENT_USER }]
  TRIGGER nome_trigger
  { BEFORE | AFTER }
  { INSERT | UPDATE | DELETE }
  ON tabella_target
  FOR EACH ROW
    trigger_stmt;
  \end{lstlisting}
  \vspace{0.2cm}
  \begin{columns}[T]
    \column{0.5\textwidth}
      \begin{itemize}
        \item \textbf{DEFINER}: account utente i cui privilegi vengono
              verificati all'attivazione
        \item \textbf{BEFORE | AFTER}: momento di esecuzione
      \end{itemize}
    \column{0.5\textwidth}
      \begin{itemize}
        \item \textbf{INSERT | UPDATE | DELETE}: evento scatenante
        \item \textbf{FOR EACH ROW}: granularità (obbligatorio in MySQL)
        \item \textbf{trigger\_stmt}: corpo del trigger
      \end{itemize}
  \end{columns}
\end{frame}

% SLIDE 20
\begin{frame}[fragile]{Parametri della Sintassi -- Dettaglio}
  \begin{description}
    \item[\texttt{DEFINER}] Determina i privilegi da applicare quando il trigger
      è attivo. Formato: \texttt{'user\_name'@'host\_name'}
    \item[\texttt{TRIGGER\_TIME}] Rappresenta il tipo di trigger:
      \texttt{BEFORE} o \texttt{AFTER}
    \item[\texttt{TRIGGER\_EVENT}] Indica il tipo di primitiva che attiva il trigger:
      \begin{itemize}
        \item \texttt{INSERT}: trigger attivato su \texttt{INSERT}, \texttt{LOAD DATA},
          \texttt{REPLACE}
        \item \texttt{UPDATE}: trigger attivato su \texttt{UPDATE}
        \item \texttt{DELETE}: trigger attivato su \texttt{DELETE}, \texttt{REPLACE}
      \end{itemize}
    \item[\texttt{TRIGGER\_STMT}] La primitiva SQL da eseguire. Per sequenze
      multiple si usa il blocco \texttt{BEGIN \ldots END}
  \end{description}
\end{frame}

% SLIDE 21
\begin{frame}[fragile]{Le Variabili OLD e NEW}
  Per accedere alle colonne della tabella target si usano gli alias \textbf{OLD}
  e \textbf{NEW}:
  \vspace{0.2cm}
  \begin{block}{\texttt{OLD.col\_name} (read-only)}
    Si riferisce alla colonna della tabella target \textbf{prima} della modifica
    effettuata dall'evento associato al trigger.
  \end{block}
  \vspace{0.2cm}
  \begin{alertblock}{\texttt{NEW.col\_name}}
    Si riferisce alla colonna della tabella target \textbf{dopo} la modifica
    effettuata dall'evento.
  \end{alertblock}
  \vspace{0.3cm}
  \begin{center}
    \renewcommand{\arraystretch}{1.2}
    \begin{tabular}{lcc}
      \toprule
      \textbf{Tipo Trigger} & \textbf{OLD disponibile} & \textbf{NEW disponibile} \\
      \midrule
      INSERT & No & Sì \\
      DELETE & Sì & No \\
      UPDATE & Sì & Sì \\
      \bottomrule
    \end{tabular}
  \end{center}
\end{frame}

% SLIDE 22
\begin{frame}[fragile]{Modifica di NEW in un BEFORE Trigger}
  \begin{block}{Caratteristica speciale dei BEFORE Trigger}
    In un \texttt{BEFORE} trigger è possibile \textbf{modificare} i valori
    che verranno scritti nella tabella usando:
    \vspace{0.2cm}
    \begin{center}
      \texttt{SET NEW.colonna = valore}
    \end{center}
  \end{block}
  \vspace{0.3cm}
  \begin{exampleblock}{Esempio: normalizzazione automatica}
    \begin{lstlisting}[style=sqlstyle]
-- Questo BEFORE INSERT trigger
-- capitalizza automaticamente il nome
CREATE TRIGGER normalizza_nome
BEFORE INSERT ON clienti
FOR EACH ROW
  SET NEW.nome = CONCAT(
    UPPER(LEFT(NEW.nome, 1)),
    LOWER(SUBSTRING(NEW.nome, 2))
  );
    \end{lstlisting}
  \end{exampleblock}
\end{frame}

% =====================================================================
\section{Privilegi e Gestione degli Errori}
% =====================================================================

% SLIDE 23
\begin{frame}{Privilegi per la Creazione di Trigger}
  \begin{columns}[T]
    \column{0.5\textwidth}
      \begin{block}{Al momento della creazione}
        L'utente che crea il trigger deve avere i \textbf{privilegi di SuperUtente}.
      \end{block}
      \vspace{0.3cm}
      \begin{block}{Clausola DEFINER}
        Specifica l'account utente associato al trigger. Formato:
        \begin{center}
          \texttt{'user'@'host'}
        \end{center}
        I privilegi dell'account DEFINER vengono verificati all'attivazione.
      \end{block}
    \column{0.5\textwidth}
      \begin{alertblock}{Privilegi necessari all'attivazione}
        L'utente DEFINER deve possedere:
        \begin{itemize}
          \item Privilegi di \textbf{SuperUtente}
          \item \textbf{SELECT} sulla tabella target (se usa OLD/NEW)
          \item \textbf{UPDATE} sulla tabella target (se usa \texttt{SET NEW.col})
          \item Tutti i privilegi per gli statement SQL nell'azione
        \end{itemize}
      \end{alertblock}
  \end{columns}
\end{frame}

% SLIDE 24
\begin{frame}{Gestione degli Errori nei Trigger}
  \begin{block}{Regole di esecuzione}
    \begin{enumerate}
      \item Se un \textbf{BEFORE trigger fallisce}, l'operazione DML che ha
            generato l'evento \textbf{non viene eseguita}.
      \item Un \textbf{AFTER trigger viene eseguito} solo se i BEFORE trigger
            e l'operazione DML hanno avuto successo.
      \item Un \textbf{errore} durante un BEFORE o AFTER trigger genera il
            fallimento dell'intera procedura.
    \end{enumerate}
  \end{block}
  \vspace{0.3cm}
  \begin{columns}[T]
    \column{0.5\textwidth}
      \begin{exampleblock}{Sistema transazionale}
        Il fallimento causa il \textbf{rollback} di \emph{tutte} le modifiche
        effettuate dal trigger.
      \end{exampleblock}
    \column{0.5\textwidth}
      \begin{alertblock}{Sistema non transazionale}
        Tutte le modifiche apportate \emph{prima} dell'errore
        \textbf{rimangono effettive}.
      \end{alertblock}
  \end{columns}
\end{frame}

% SLIDE 25
\begin{frame}{Limitazioni dei Trigger in MySQL}
  \begin{alertblock}{Unicità per evento}
    In MySQL \textbf{non possono essere creati due trigger} con evento e tipo
    evento uguali sulla stessa tabella target.
  \end{alertblock}
  \vspace{0.3cm}
  \begin{exampleblock}{Esempi di combinazioni non ammesse}
    Non possono coesistere sulla stessa tabella:
    \begin{itemize}
      \item Due trigger \texttt{BEFORE INSERT}
      \item Due trigger \texttt{AFTER UPDATE}
      \item Due trigger \texttt{BEFORE DELETE}
    \end{itemize}
  \end{exampleblock}
  \vspace{0.3cm}
  \begin{block}{Altre limitazioni}
    \begin{itemize}
      \item I trigger \textbf{non possono invocare} procedure che restituiscono
            dati all'utente o che usano SQL dinamico
      \item I trigger \textbf{non possono invocare} procedure che usano
            \texttt{START TRANSACTION}, \texttt{COMMIT} o \texttt{ROLLBACK}
    \end{itemize}
  \end{block}
\end{frame}

% SLIDE 26
\begin{frame}{Chiamata di Procedure nei Trigger}
  \begin{columns}[T]
    \column{0.5\textwidth}
      \begin{alertblock}{Procedure NON invocabili}
        \begin{itemize}
          \item Procedure che \textbf{restituiscono dati} all'utente
          \item Procedure che usano \textbf{SQL dinamico}
          \item Procedure con \texttt{START TRANSACTION}
          \item Procedure con \texttt{END TRANSACTION}
          \item Procedure con \texttt{COMMIT}
          \item Procedure con \texttt{ROLLBACK}
        \end{itemize}
      \end{alertblock}
    \column{0.5\textwidth}
      \begin{exampleblock}{Costrutti utilizzabili}
        All'interno del corpo del trigger è possibile utilizzare:
        \begin{itemize}
          \item Blocco \texttt{BEGIN \ldots END}
          \item \texttt{IF \ldots THEN \ldots ELSE}
          \item \texttt{DECLARE} per variabili locali
          \item \texttt{SELECT \ldots INTO}
          \item \texttt{INSERT}, \texttt{UPDATE}, \texttt{DELETE}
          \item Chiamata a \textbf{funzioni} definite dall'utente
        \end{itemize}
      \end{exampleblock}
  \end{columns}
\end{frame}

% =====================================================================
\section{Esempi Pratici}
% =====================================================================

% SLIDE 27
\begin{frame}[fragile]{Esempio Semplice: Trigger su INSERT}
  \begin{lstlisting}[style=sqlstyle]
-- Creazione tabella
CREATE TABLE account (
  acct_num INT,
  amount   DECIMAL(10,2)
);

-- Trigger: accumula la somma degli importi inseriti
CREATE TRIGGER ins_sum
BEFORE INSERT ON account
FOR EACH ROW
  SET @sum = @sum + NEW.amount;
  \end{lstlisting}
  \begin{block}{Utilizzo}
    \begin{lstlisting}[style=sqlstyle]
SET @sum = 0;
INSERT INTO account VALUES (1, 14.98), (2, 1937.50);
SELECT @sum;  -- Risultato: 1952.48
    \end{lstlisting}
  \end{block}
\end{frame}

% SLIDE 28
\begin{frame}[fragile]{Esempio 1 -- Tabella Utente}
  \begin{lstlisting}[style=sqlstyle]
CREATE TABLE `bd2`.`utente` (
  `idutente`  INTEGER UNSIGNED NOT NULL AUTO_INCREMENT,
  `nome`      VARCHAR(15) NOT NULL DEFAULT '',
  `cognome`   VARCHAR(25) NOT NULL DEFAULT '',
  `password`  VARCHAR(45) NOT NULL DEFAULT '',
  PRIMARY KEY (`idutente`)
) ENGINE = InnoDB;
  \end{lstlisting}
  \begin{block}{Obiettivo del trigger}
    Ogni volta che viene inserito un utente:
    \begin{enumerate}
      \item Il nome e il cognome vengono \textbf{capitalizzati} automaticamente
      \item La password viene \textbf{hashata} con la funzione \texttt{PASSWORD()}
    \end{enumerate}
  \end{block}
\end{frame}

% SLIDE 29
\begin{frame}[fragile]{Esempio 1 -- Funzione di Validazione}
  \begin{lstlisting}[style=sqlstyle]
DELIMITER //
CREATE FUNCTION valida_stringa (str VARCHAR(25))
RETURNS VARCHAR(25)
DETERMINISTIC
BEGIN
  DECLARE testa VARCHAR(1);
  DECLARE coda  VARCHAR(24);

  SET str   = TRIM(str);
  SET testa = UPPER(LEFT(str, 1));
  SET coda  = LOWER(RIGHT(str, LENGTH(str) - 1));

  RETURN CONCAT(testa, coda);
END;
//
DELIMITER ;
  \end{lstlisting}
  \begin{block}{Comportamento}
    La funzione trasforma qualunque stringa in formato \textbf{Title Case}:
    prima lettera maiuscola, resto minuscolo, spazi eliminati.
  \end{block}
\end{frame}

% SLIDE 30
\begin{frame}[fragile]{Esempio 1 -- Creazione del Trigger}
  \begin{lstlisting}[style=sqlstyle]
DELIMITER //
CREATE TRIGGER UTENTE_BEF_INS_ROW
BEFORE INSERT ON utente
FOR EACH ROW
BEGIN
  SET NEW.nome     = valida_stringa(NEW.nome);
  SET NEW.cognome  = valida_stringa(NEW.cognome);
  SET NEW.password = PASSWORD(NEW.password);
END;
//
DELIMITER ;
  \end{lstlisting}
  \begin{alertblock}{Nota sulla convenzione di naming}
    Il nome \texttt{UTENTE\_BEF\_INS\_ROW} codifica:
    tabella (\texttt{UTENTE}), momento (\texttt{BEF}),
    evento (\texttt{INS}), granularità (\texttt{ROW}).
  \end{alertblock}
\end{frame}

% SLIDE 31
\begin{frame}[fragile]{Esempio 1 -- Inserimento e Risultato}
  \begin{lstlisting}[style=sqlstyle]
-- Inserimento con dati non normalizzati
INSERT INTO utente VALUES (1, 'alessio',    'meola',       'passwd1010');
INSERT INTO utente VALUES (2, 'gIuSePPe',   'MAZZARELLA',  'passwd2020');
  \end{lstlisting}
  \vspace{0.2cm}
  \begin{exampleblock}{Risultato dopo SELECT * FROM utente}
    \begin{lstlisting}[style=outputstyle]
+----------+----------+------------+------------------------------------+
| idutente | nome     | cognome    | password                           |
+----------+----------+------------+------------------------------------+
|        1 | Alessio  | Meola      | *4BD714D8E708D3D6BF85BB246E0DCE6C32|
|        2 | Giuseppe | Mazzarella | *80BCC99B46B82A5A8C18B98EDD7585AE51|
+----------+----------+------------+------------------------------------+
    \end{lstlisting}
  \end{exampleblock}
  \begin{block}{Osservazioni}
    Il trigger ha automaticamente capitalizzato i nomi e hashato le password,
    indipendentemente dal formato di input fornito.
  \end{block}
\end{frame}

% SLIDE 32
\begin{frame}{Esempio 2 -- Gestione Magazzino}
  \begin{block}{Scenario}
    Un sistema di gestione magazzino deve gestire automaticamente:
    \begin{enumerate}
      \item \textbf{Riordino della merce}: quando la quantità scende sotto
            il limite minimo, viene creato automaticamente un ordine
      \item \textbf{Statistiche sulle vendite}: aggiornamento automatico
            delle analisi per prodotto
    \end{enumerate}
  \end{block}
  \vspace{0.2cm}
  \begin{alertblock}{Soluzione implementata}
    \begin{itemize}
      \item \textbf{2 trigger}: \texttt{PROD\_MAG\_AFT\_UPD\_ROW} e
            \texttt{PROD\_MAG\_AFT\_INS\_ROW}
      \item \textbf{1 funzione}: \texttt{analisi\_vendite()}
    \end{itemize}
  \end{alertblock}
\end{frame}

% SLIDE 33
\begin{frame}[fragile]{Esempio 2 -- Schema delle Tabelle}
  \begin{lstlisting}[style=sqlstyle]
-- Prodotti disponibili
CREATE TABLE prodotto (
  idprodotto INTEGER UNSIGNED NOT NULL AUTO_INCREMENT,
  nome       VARCHAR(45) NOT NULL,
  marca      VARCHAR(45) NOT NULL,
  vendita    FLOAT       NOT NULL DEFAULT 0,  -- prezzo vendita
  acquisto   FLOAT       NOT NULL DEFAULT 0,  -- prezzo acquisto
  PRIMARY KEY(idprodotto)
) ENGINE = InnoDB;

-- Depositi/Magazzini
CREATE TABLE magazzino (
  idmagazzino INTEGER UNSIGNED NOT NULL AUTO_INCREMENT,
  citta       VARCHAR(45) NOT NULL,
  PRIMARY KEY(idmagazzino)
) ENGINE = InnoDB;
  \end{lstlisting}
\end{frame}

% SLIDE 34
\begin{frame}[fragile]{Esempio 2 -- Tabella Prodotti\_Magazzino}
  \begin{lstlisting}[style=sqlstyle]
CREATE TABLE prodotti_magazzino (
  idprodotto   INTEGER UNSIGNED NOT NULL AUTO_INCREMENT,
  idmagazzino  INTEGER UNSIGNED NOT NULL DEFAULT 0,
  qta          INTEGER UNSIGNED NOT NULL DEFAULT 0,
  qta_limite   INTEGER UNSIGNED NOT NULL DEFAULT 0,
  qta_riordino INTEGER UNSIGNED NOT NULL DEFAULT 0,
  PRIMARY KEY(idprodotto, idmagazzino),
  CONSTRAINT prod_mag_prod FOREIGN KEY (idprodotto)
    REFERENCES prodotto(idprodotto)
    ON DELETE RESTRICT ON UPDATE RESTRICT,
  CONSTRAINT prod_mag_mag FOREIGN KEY (idmagazzino)
    REFERENCES magazzino(idmagazzino)
    ON DELETE RESTRICT ON UPDATE RESTRICT
) ENGINE = InnoDB;
  \end{lstlisting}
  \begin{block}{Chiave}
    \texttt{qta\_limite}: soglia minima sotto cui scatta il riordino.\\
    \texttt{qta\_riordino}: quantità da ordinare al fornitore.
  \end{block}
\end{frame}

% SLIDE 35
\begin{frame}[fragile]{Esempio 2 -- Tabelle Ordine e Analisi}
  \begin{lstlisting}[style=sqlstyle]
-- Ordini di rifornitamento generati automaticamente
CREATE TABLE ordine (
  idprodotto  INTEGER UNSIGNED NOT NULL AUTO_INCREMENT,
  idmagazzino INTEGER UNSIGNED NOT NULL DEFAULT 0,
  data        DATETIME         NOT NULL DEFAULT 0,
  PRIMARY KEY(idmagazzino, idprodotto),
  CONSTRAINT ord_prod FOREIGN KEY (idprodotto)
    REFERENCES prodotto(idprodotto),
  CONSTRAINT ord_mag FOREIGN KEY (idmagazzino)
    REFERENCES magazzino(idmagazzino)
) ENGINE = InnoDB;

-- Statistiche di vendita per prodotto
CREATE TABLE analisi (
  idprodotto INTEGER UNSIGNED NOT NULL AUTO_INCREMENT,
  venduto    INTEGER UNSIGNED NOT NULL DEFAULT 0,
  guadagno   FLOAT            NOT NULL DEFAULT 0,
  PRIMARY KEY(idprodotto),
  CONSTRAINT analisi_prod FOREIGN KEY (idprodotto)
    REFERENCES prodotto(idprodotto)
) ENGINE = InnoDB;
  \end{lstlisting}
\end{frame}

% SLIDE 36
\begin{frame}[fragile]{Esempio 2 -- Funzione analisi\_vendite()}
  \begin{lstlisting}[style=sqlstyle]
DELIMITER //
CREATE FUNCTION analisi_vendite(prodotto INT, qta INT)
RETURNS INTEGER
DETERMINISTIC
BEGIN
  DECLARE guad INTEGER;

  -- Calcola il guadagno unitario (vendita - acquisto)
  SELECT (vendita - acquisto) INTO guad
  FROM prodotto WHERE idprodotto = prodotto;

  -- Aggiorna le statistiche
  UPDATE analisi
  SET venduto  = venduto + qta,
      guadagno = guadagno + (guad * qta)
  WHERE idprodotto = prodotto;

  RETURN 1;
END;
//
DELIMITER ;
  \end{lstlisting}
\end{frame}

% SLIDE 37
\begin{frame}[fragile]{Esempio 2 -- Trigger Riordino (AFTER UPDATE)}
  \begin{lstlisting}[style=sqlstyle]
DELIMITER //
CREATE TRIGGER PROD_MAG_AFT_UPD_ROW
AFTER UPDATE OF qta ON prodotti_magazzino
FOR EACH ROW
BEGIN
  DECLARE nordini    INTEGER;
  DECLARE qtavenduta INTEGER;

  SET qtavenduta = OLD.qta - NEW.qta;  -- quantita' venduta

  -- Controlla se la scorta e' sotto il limite
  IF NEW.qta < NEW.qta_limite THEN
    SELECT COUNT(*) INTO nordini FROM ordine
    WHERE ordine.idprodotto  = NEW.idprodotto
      AND ordine.idmagazzino = NEW.idmagazzino;

    -- Crea ordine solo se non ne esiste gia' uno
    IF nordini = 0 THEN
      INSERT INTO ordine VALUES
        (NEW.idprodotto, NEW.idmagazzino, CURDATE());
    END IF;
  END IF;

  CALL analisi_vendite(NEW.idprodotto, qtavenduta);
END;
//
  \end{lstlisting}
\end{frame}

% SLIDE 38
\begin{frame}[fragile]{Esempio 2 -- Trigger Inizializzazione Analisi}
  \begin{lstlisting}[style=sqlstyle]
-- Trigger: inizializza il record in "analisi"
-- quando un prodotto viene aggiunto al magazzino
CREATE TRIGGER PROD_MAG_AFT_INS_ROW
AFTER INSERT ON prodotti_magazzino
FOR EACH ROW
BEGIN
  IF NOT EXISTS (
    SELECT * FROM analisi
    WHERE idprodotto = NEW.idprodotto
  ) THEN
    INSERT INTO analisi VALUES (NEW.idprodotto, 0, 0);
  END IF;
END;
//
DELIMITER ;
  \end{lstlisting}
  \begin{block}{Logica}
    Quando un prodotto viene inserito in magazzino per la prima volta,
    viene creato automaticamente un record iniziale nella tabella \texttt{analisi}
    con quantità venduta e guadagno pari a zero.
  \end{block}
\end{frame}

% SLIDE 39
\begin{frame}[fragile]{Esempio 2 -- Esecuzione e Risultati}
  \begin{lstlisting}[style=sqlstyle]
-- Vendita: riduzione scorte prodotto 3, magazzino 1
UPDATE prodotti_magazzino
SET qta = qta - 41
WHERE idprodotto = 3 AND idmagazzino = 1;
  \end{lstlisting}
  \begin{columns}[T]
    \column{0.5\textwidth}
      \begin{exampleblock}{Tabella ordine (auto-generato)}
        \begin{lstlisting}[style=outputstyle]
+----+-----+------------+
|prod|mag  |data        |
+----+-----+------------+
|  3 |  1  |2006-01-13  |
|  1 |  2  |2006-01-13  |
+----+-----+------------+
        \end{lstlisting}
      \end{exampleblock}
    \column{0.5\textwidth}
      \begin{exampleblock}{Tabella analisi (aggiornata)}
        \begin{lstlisting}[style=outputstyle]
+----+--------+---------+
|prod|venduto |guadagno |
+----+--------+---------+
|  1 |      8 |     400 |
|  2 |      0 |       0 |
|  3 |     41 |    3280 |
+----+--------+---------+
        \end{lstlisting}
      \end{exampleblock}
  \end{columns}
\end{frame}

% =====================================================================
\section{Riepilogo}
% =====================================================================

% SLIDE 40
\begin{frame}{Riepilogo -- Trigger in MySQL}
  \begin{columns}[T]
    \column{0.5\textwidth}
      \textbf{Punti chiave}
      \begin{itemize}
        \item I trigger implementano il paradigma \textbf{E-C-A}
        \item Le basi di dati con trigger sono dette \textbf{attive}
        \item MySQL supporta solo trigger \textbf{row-level}
        \item Le variabili \textbf{OLD} e \textbf{NEW} accedono ai dati
        \item \textbf{BEFORE}: modifica dati prima della scrittura
        \item \textbf{AFTER}: azioni collaterali post-scrittura
      \end{itemize}
    \column{0.5\textwidth}
      \textbf{Best Practice}
      \begin{itemize}
        \item Naming convention chiara (tabella\_BEF/AFT\_INS/UPD/DEL\_ROW)
        \item Evitare logiche complesse nei trigger
        \item Attenzione ai trigger in cascata
        \item Testare il comportamento su errori
        \item Un solo trigger per combinazione evento/tempo/tabella
      \end{itemize}
  \end{columns}
  \vspace{0.4cm}
  \begin{center}
    \begin{tikzpicture}
      \node[draw,fill=sqlblue!20,rounded corners,minimum width=2cm] (e) at (0,0) {Evento};
      \node[draw,fill=orange!20,rounded corners,minimum width=2cm] (c) at (3,0) {Condizione};
      \node[draw,fill=sqlgreen!20,rounded corners,minimum width=2cm] (a) at (6,0) {Azione};
      \draw[->,thick] (e)--(c);
      \draw[->,thick] (c)--(a);
    \end{tikzpicture}
  \end{center}
\end{frame}

\end{document}
